
\subsubsection{5.1 Alignment Increases Brier Reliability (H-E1)}

Eight of nine aligned model-size pairs showed positive $\Delta Brier$ Reliability, indicating that alignment training consistently increases overconfidence relative to paired base models on MMLU. The exception was 6.9B-PPO ($\Delta Reliability$ = -0.0036), which showed a marginal calibration improvement.

\textbf{Table 1.} $\Delta ECE$ and $\Delta Brier$ Reliability for all nine alignment-size pairs on MMLU.

\begin{table*}[htbp]
\centering
\resizebox{\dimexpr 2\columnwidth+\columnsep\relax}{!}{%
\begin{tabular}{lllllll}
\toprule
Model & Alignment & ECE\_base & ECE\_aligned & $\Delta ECE$ & $\Delta Brier$ Reliability & 95\% CI Lower \\
\midrule
Pythia-1.4B & SFT & 0.0849 & 0.1415 & +0.0566 & +0.0289 & +0.0269 \\
Pythia-1.4B & DPO & 0.0849 & 0.2516 & +0.1667 & +0.1048 & +0.1009 \\
Pythia-1.4B & PPO & 0.0849 & 0.1923 & +0.1074 & +0.0406 & +0.0345 \\
Pythia-2.8B & SFT & 0.0597 & 0.0694 & +0.0097 & +0.0033 & +0.0021 \\
Pythia-2.8B & DPO & 0.0597 & 0.1441 & +0.0844 & +0.0437 & +0.0407 \\
Pythia-2.8B & PPO & 0.0597 & 0.1577 & +0.0980 & +0.0423 & +0.0388 \\
Pythia-6.9B & SFT & 0.0792 & 0.0830 & +0.0038 & +0.0010 & +0.0001 \\
Pythia-6.9B & DPO & 0.0792 & 0.1010 & +0.0218 & +0.0099 & +0.0090 \\
Pythia-6.9B & PPO & 0.0792 & 0.0609 & -0.0183 & -0.0036 & -0.0053 \\
\bottomrule
\end{tabular}
}
\end{table*}

The H-E1 MUST\_WORK gate (bootstrap CI lower > 0 for PPO or DPO in >= 2/3 sizes) was satisfied via both DPO (3/3 sizes) and PPO (2/3 sizes).

\textbf{Causal baseline confirmation (H-M1):} Base Pythia models showed ECE = 0.0849 (1.4B), 0.0597 (2.8B), 0.0792 (6.9B), all below the 0.15 threshold (MUST\_WORK PASS). Pretraining yielded well-calibrated logits; the observed increases are attributable to alignment.

\begin{figure}[htbp]
\centering
\includegraphics[width=\columnwidth]{figures/delta_reliability_bar.png}
\caption{Figure 1: $\Delta Brier$ Reliability across all alignment-size pairs}
\label{fig:delta_reliability_bar}
\end{figure}

\textit{Figure 1. $\Delta Brier$ Reliability (aligned minus base) for all nine alignment-size pairs on MMLU. Error bars show bootstrap 95\% confidence intervals (n = 1,000). Positive values indicate increased overconfidence.}

\begin{figure}[htbp]
\centering
\includegraphics[width=\columnwidth]{figures/brier_decomposition.png}
\caption{Figure 2: Brier score decomposition}
\label{fig:brier_decomposition}
\end{figure}

\textit{Figure 2. Brier score decomposition (Reliability, Resolution, Uncertainty) for all 12 models. Reliability increases under alignment while Resolution changes are moderate, indicating that overconfidence rather than discriminability collapse is the primary calibration change.}

\subsubsection{5.2 Exploratory Ordering: DPO >= PPO > SFT (H-M2)}

The pre-registered calibration ordering prediction (PPO >= DPO > SFT, based on conventional reward optimization pressure) was not supported by the data. DPO produced larger calibration degradation than PPO in all three model sizes as measured by $\Delta Brier$ Reliability, yielding an observed ordering of DPO >= PPO > SFT.

Pre-softmax logit margin analysis confirmed that confidence inflation operates at the logit level, not as a softmax normalization artifact. DPO showed positive $\Delta margin$ in all three sizes; PPO in two of three sizes (6.9B-PPO showed $\Delta margin$ = -0.036).

\begin{figure}[htbp]
\centering
\includegraphics[width=\columnwidth]{figures/figure_01_delta_margin_gate.png}
\caption{Figure 3: Pre-softmax margin inflation}
\label{fig:figure_01_delta_margin_gate}
\end{figure}

\textit{Figure 3. Pre-softmax logit margin inflation ($\Delta margin$ = margin\_aligned - margin\_base) with bootstrap 95\% confidence intervals. DPO shows positive $\Delta margin$ in all three sizes; PPO in two of three sizes.}

\begin{figure}[htbp]
\centering
\includegraphics[width=\columnwidth]{figures/figure_04_gradient_ordering_heatmap.png}
\caption{Figure 4: Gradient ordering heatmap}
\label{fig:figure_04_gradient_ordering_heatmap}
\end{figure}

\textit{Figure 4. $\Delta margin$ heatmap by alignment method (rows) and Pythia model size (columns).}

\textbf{Statistical qualification.} At 2.8B, the DPO-PPO difference was narrow: $\Delta Rel_DPO$ = 0.0437 (95\% CI: [0.0407, 0.0469]) versus $\Delta Rel_PPO$ = 0.0423 (95\% CI: [0.0388, 0.0456]). These confidence intervals overlap, and the difference of 0.0014 is not statistically distinguishable at the 95\% confidence level. The DPO >= PPO ordering is directionally consistent across all three sizes but should be treated as inconclusive at 2.8B. The finding is primarily supported by the clear separation at 1.4B (DPO: 0.1048, PPO: 0.0406) and 6.9B (DPO: +0.0099, PPO: -0.0036).

The margin analysis (H-M2) and Spearman rho analysis (H-M3) are both derived from the same per-item log-probability vectors and should be understood as complementary characterizations of the same logit distribution changes rather than independent confirmatory evidence.

\subsubsection{5.3 Mechanism Discrimination: H2 Dominance (H-M3)}

All nine alignment-size pairs fell below the H1 threshold (rho >= 0.90). Eight of nine fell below the H2 diagnostic threshold (rho < 0.85).

\textbf{Table 2.} H1/H2 discrimination summary for all nine alignment-size pairs.

\begin{table*}[htbp]
\centering
\resizebox{\dimexpr 2\columnwidth+\columnsep\relax}{!}{%
\begin{tabular}{lllll}
\toprule
Pair & Spearman rho & H1 (>= 0.90) & H2 (< 0.85) & Argmax Changed (\%) \\
\midrule
1.4B-SFT & 0.753 & No & Yes & 42.8\% (6,014 / 14,042) \\
1.4B-DPO & 0.737 & No & Yes & 41.4\% (5,807 / 14,042) \\
1.4B-PPO & -0.324 & No & Yes & 99.7\% (13,998 / 14,042) \\
2.8B-SFT & 0.719 & No & Yes & 28.7\% (4,025 / 14,042) \\
2.8B-DPO & 0.590 & No & Yes & 39.4\% (5,526 / 14,042) \\
2.8B-PPO & 0.175 & No & Yes & 64.5\% (9,049 / 14,042) \\
6.9B-SFT & 0.839 & No & Yes & 14.8\% (2,073 / 14,042) \\
6.9B-DPO & 0.875 & No & No (near threshold) & 15.8\% (2,224 / 14,042) \\
6.9B-PPO & 0.505 & No & Yes & 39.7\% (5,570 / 14,042) \\
\bottomrule
\end{tabular}
}
\end{table*}

The 1.4B-PPO case is particularly notable: rho = -0.324 indicates that the aligned model's answer preferences are negatively correlated with the base model's. Only 44 of 14,042 MMLU items (0.3\%) maintain the same top-ranked answer option after PPO alignment at 1.4B. The 6.9B-DPO model showed the highest rho (0.875), the closest to the H1 threshold, suggesting a possible scale-mediated effect.

\begin{figure}[htbp]
\centering
\includegraphics[width=\columnwidth]{figures/figure_01_spearman_rho.png}
\caption{Figure 5: Spearman rho per alignment-size pair}
\label{fig:figure_01_spearman_rho}
\end{figure}

\textit{Figure 5. Spearman rank correlation (rho) between base and aligned 4-option log-probability vectors per MMLU item. Dashed lines at rho = 0.90 (H1 threshold) and rho = 0.85 (H2 diagnostic). All nine pairs fall below 0.90; eight of nine fall below 0.85.}

\begin{figure}[htbp]
\centering
\includegraphics[width=\columnwidth]{figures/figure_04_argmax_proportion.png}
\caption{Figure 6: Argmax redistribution rates}
\label{fig:figure_04_argmax_proportion}
\end{figure}

\textit{Figure 6. Proportion of MMLU items where alignment changes the argmax prediction. 1.4B-PPO changes argmax for 99.7\% of items.}

\begin{figure}[htbp]
\centering
\includegraphics[width=\columnwidth]{figures/figure_03_brier_partition.png}
\caption{Figure 7: Brier partition by argmax agreement}
\label{fig:figure_03_brier_partition}
\end{figure}

\textit{Figure 7. Brier reliability partitioned into shared-argmax and changed-argmax items. For 1.4B-PPO, the shared-argmax partition contains only 44 items.}

\subsubsection{5.4 H3 Exclusion in Softmax-ECE Setting}

$\Delta ECE$ on TruthfulQA MC1 was less than $\Delta ECE$ on MMLU for all three alignment types (SFT ratio = 0.32, DPO ratio = 0.26, PPO ratio = 0.73). The H3 hypothesis predicts ratios >= 1.0 if framing susceptibility is the primary driver. The observed ratios are inconsistent with H3 as a dominant mechanism in the softmax-ECE evaluation setting.

This result is distinct from the findings of Chhikara et al. (2025), who observed framing sensitivity using verbally elicited confidence. The measurement modality (softmax log-probability ECE versus verbal confidence) appears to determine whether framing effects are observed. A methodological caveat applies: the shot-count asymmetry between MMLU (4-shot) and TruthfulQA MC1 (0-shot) could contribute to the $\Delta ECE_MMLU$ > $\Delta ECE_TruthfulQA$ pattern independently of framing susceptibility, though the magnitude of the ratios (0.26-0.73) suggests framing is unlikely to be the dominant driver even accounting for this difference.

\begin{figure}[htbp]
\centering
\includegraphics[width=\columnwidth]{figures/figure_05_truthfulqa_ece.png}
\caption{Figure 8: TruthfulQA H3 diagnostic}
\label{fig:figure_05_truthfulqa_ece}
\end{figure}

\textit{Figure 8. $\Delta ECE$ on TruthfulQA MC1 versus MMLU for all alignment types. $\Delta ECE_TruthfulQA$ < $\Delta ECE_MMLU$ for all three methods.}
